%!TEX program = pdflatex
% Compile in Kile with -shell-escape (minted) and biber; present with pdfpc.
% Movies follow the deck convention: an MP4 and a same-basename PNG poster in
% Movies/, embedded as a run: link annotation that pdfpc turns into the video.
\documentclass[aspectratio=169]{beamer}

\usetheme{Janelia}
\usepackage{people}

\usepackage[english]{babel}
\usepackage[T1]{fontenc}
\usepackage[overlay,absolute]{textpos}
\setlength{\TPHorizModule}{1cm}
\setlength{\TPVertModule}{1cm}

\usepackage{amsmath,amssymb}
\usepackage{multimedia}
\usepackage{tabularx,booktabs}
\usepackage{helvet}
\usepackage{tikz}
\usetikzlibrary{arrows.meta,positioning}

\newcommand{\playmovie}[3][\textwidth]{%
  \href{run:#2.mp4?autostart&loop&noprogress}{\includegraphics[width=#1]{#2.png}}%
  \\[2pt]{\scriptsize #3}}

\title{Recovering synaptic conductances from dynamics}
\subtitle{A conductance GNN on task-trained \emph{Drosophila} visual circuitry}
\author{Cedric Allier}
\institute{Saalfeld lab, Janelia}
\date{\today}

\begin{document}

\frame{\titlepage}

% ===================================================================
\section{Generating conductance training data}
% ===================================================================

\begin{frame}{A conductance flyvis model, and how its data is made}
\begin{columns}[T]
\begin{column}{0.52\textwidth}
{\small\textbf{Nominal flyvis} --- current-based synapse:}
\[ \tau_i \dot v_i = -(v_i - V^{\mathrm{rest}}_i) + \sum_j W_{ij}\,\mathrm{relu}(v_j) + I_i(t) \]
{\small\textbf{This model} --- conductance synapse, $W_{ij}\ge 0$:}
\[ \tau_i \dot v_i = -(v_i - V^{\mathrm{rest}}_i) + \sum_j W_{ij}\,\mathrm{relu}(v_j)\,(E_i - v_i) + I_i(t) \]
{\small The polarity of a synapse moves out of $W$ and into the driving force
$(E_i-v_i)$: $W$ becomes a non-negative conductance and the sign is carried by
the postsynaptic reversal potential.}
\end{column}
\begin{column}{0.46\textwidth}
{\small\textbf{Integration: exponential Euler.}\\[3pt]
Collecting the terms in $v_i$,
\[ \tau_i \dot v_i = -(1+G_i)\,v_i + b_i, \qquad G_i=\textstyle\sum_j W_{ij}\mathrm{relu}(v_j) \]
is linear at frozen coefficients, so the step is exact:
\[ v_i \leftarrow v_i^{\infty} + (v_i - v_i^{\infty})e^{-z_i},\quad z_i=\tfrac{\Delta t}{\tau_i}(1+G_i) \]
\textbf{Why it matters.} Forward Euler contracts only while $z<2$; this network
reaches $z=4.42$. A model that learned the generator exactly would still diverge
in a forward-Euler rollout at $\Delta t = 20$\,ms.}
\end{column}
\end{columns}
\end{frame}

\begin{frame}{Reversal potentials by cell type}
\begin{center}
% TODO asset: E_i vs postsynaptic type for the flowcond generator.
% reversals.png exists only for flyvis_conductance_ion_sub_*, a DIFFERENT
% dataset (the teacher-student twin) -- do not substitute it here.
% The values live in graphs_data/fly/flyvis_flowcond_*/ode_params.pt.
\fbox{\parbox[c][4.5cm][c]{0.7\textwidth}{\centering
  \textbf{asset missing}\\[4pt]
  \texttt{reversals.png} has never been written for the \texttt{flowcond}
  series --- only for the \texttt{ion\_sub} twin.\\
  $E_i$ per type is in \texttt{ode\_params.pt}.}}
\end{center}
\end{frame}

\begin{frame}{Activity: noise free and $\sigma = 0.05$}
\begin{columns}[T]
\begin{column}{0.5\textwidth}
\centering\includegraphics[width=\textwidth]{Fig/activity_noisefree.png}\\[2pt]
{\scriptsize noise free}
\end{column}
\begin{column}{0.5\textwidth}
\centering\includegraphics[width=\textwidth]{Fig/activity_noise005.png}\\[2pt]
{\scriptsize $\sigma = 0.05$ process noise}
\end{column}
\end{columns}
\vspace{4pt}
{\scriptsize The noise is added to the voltage at every integration step, so the
stored trajectory is one stochastic path. It lives in the TRAIN split only: the
test split is identical to the noise-free twin to $2.6\times10^{-6}$.}
\end{frame}

\begin{frame}{Generation}
\begin{center}
\playmovie[0.62\textwidth]{Movies/generation_flowcond}{%
  \texttt{flyvis\_flowcond\_noise\_005\_blank50\_cv00} --- 13{,}741 neurons,
  434{,}112 edges, 64{,}000 frames at $\Delta t = 20$\,ms, DAVIS video with a
  50\% blank prefix.}
\end{center}
\end{frame}

% ===================================================================
\section{Conductance model on current data}
% ===================================================================

\begin{frame}{Conductance model fitted to current-generated data}
% TODO: 1-3 slides. Source: config/fly/archive_7/flyvis_current_*_conductance_*.
\end{frame}

% ===================================================================
\section{Known-ODE on conductance data}
% ===================================================================

\begin{frame}{The identifiability ceiling}
% TODO: 10 folds RUNNING as of 2026-09-21
% flyvis_flowcond_{noise_005,noise_free}_conductance_knownode_cv0{0..4}.
% The known-ODE has the generator's structure and learns only its parameters,
% so it separates the inverse problem from the readout.
\end{frame}

% ===================================================================
\section{Conductance GNN on conductance data}
% ===================================================================

\begin{frame}{What the GNN recovers}
% TODO: Wij_R2 peak 0.433 +- 0.074 (n=22 seeds), msg_i_R2 0.90-0.94,
% rollout_r 0.99 and saturated. The sum of a neuron's messages is recovered;
% its decomposition into edges is not.
\end{frame}

\begin{frame}{Two failure modes}
% TODO: sign-flip basin ~41% of seeds (k_i negative, wrong-slope 65-68%,
% |W| and hence Wij_R2 unharmed); W-scale overshoot (gain -> 1, tau_R2 0.30).
\end{frame}

% ===================================================================
\section{Perspectives}
% ===================================================================

\begin{frame}{Where the ceiling comes from, and what would move it}
% TODO: separable g_phi = phi(v_j)*psi(v_i); decorrelating stimuli;
% presynaptic silencing in the generator.
\end{frame}

\end{document}
